% Sections won't have numbers
% Make the page area as large as possible
% Set the first level numbering to arabic and the second level
% to alphabetic. The remaining two levels are arabic
%Change to \pagestyle{empty} to suppress numbers on following pages
%\pagestyle{empty}
% Do not reset outermost enumerate counter


\documentclass[onecolumn]{article}
%%%%%%%%%%%%%%%%%%%%%%%%%%%%%%%%%%%%%%%%%%%%%%%%%%%%%%%%%%%%%%%%%%%%%%%%%%%%%%%%%%%%%%%%%%%%%%%%%%%%%%%%%%%%%%%%%%%%%%%%%%%%%%%%%%%%%%%%%%%%%%%%%%%%%%%%%%%%%%%%%%%%%%%%%%%%%%%%%%%%%%%%%%%%%%%%%%%%%%%%%%%%%%%%%%%%%%%%%%%%%%%%%%%%%%%%%%%%%%%%%%%%%%%%%%%%
\usepackage{amsmath}
\usepackage{sw20exm2}
\usepackage[compat2]{geometry}

\setcounter{MaxMatrixCols}{10}
%TCIDATA{OutputFilter=LATEX.DLL}
%TCIDATA{Version=5.50.0.2953}
%TCIDATA{<META NAME="SaveForMode" CONTENT="1">}
%TCIDATA{BibliographyScheme=Manual}
%TCIDATA{Created=Sunday, August 26, 2012 12:26:36}
%TCIDATA{LastRevised=Wednesday, April 01, 2026 13:45:38}
%TCIDATA{<META NAME="ViewSettings" CONTENT="0">}
%TCIDATA{<META NAME="GraphicsSave" CONTENT="32">}
%TCIDATA{<META NAME="DocumentShell" CONTENT="Exams and Syllabi\SW\SW Exam #1 - 8.5x14 Page">}
%TCIDATA{CSTFile=Exam.cst}

\setlength{\topmargin}{-0.75in}
\setlength{\textheight}{12.25in}
\setlength{\oddsidemargin}{0.0in}
\setlength{\evensidemargin}{0.0in}
\setlength{\textwidth}{6.5in}
\def\labelenumi{\arabic{enumi}.}
\def\theenumi{\arabic{enumi}}
\def\labelenumii{(\alph{enumii})}
\def\theenumii{\alph{enumii}}
\def\p@enumii{\theenumi.}
\def\labelenumiii{\arabic{enumiii}.}
\def\theenumiii{\arabic{enumiii}}
\def\p@enumiii{(\theenumi)(\theenumii)}
\def\labelenumiv{\arabic{enumiv}.}
\def\theenumiv{\arabic{enumiv}}
\def\p@enumiv{\p@enumiii.\theenumiii}
\pagestyle{plain}
\setcounter{secnumdepth}{0}
\newtheorem{theorem}{Theorem}
\newtheorem{acknowledgement}[theorem]{Acknowledgement}
\newtheorem{algorithm}[theorem]{Algorithm}
\newtheorem{axiom}[theorem]{Axiom}
\newtheorem{case}[theorem]{Case}
\newtheorem{claim}[theorem]{Claim}
\newtheorem{conclusion}[theorem]{Conclusion}
\newtheorem{condition}[theorem]{Condition}
\newtheorem{conjecture}[theorem]{Conjecture}
\newtheorem{corollary}[theorem]{Corollary}
\newtheorem{criterion}[theorem]{Criterion}
\newtheorem{definition}[theorem]{Definition}
\newtheorem{example}[theorem]{Example}
\newtheorem{exercise}[theorem]{Exercise}
\newtheorem{lemma}[theorem]{Lemma}
\newtheorem{notation}[theorem]{Notation}
\newtheorem{problem}[theorem]{Problem}
\newtheorem{proposition}[theorem]{Proposition}
\newtheorem{remark}[theorem]{Remark}
\newtheorem{solution}[theorem]{Solution}
\newtheorem{summary}[theorem]{Summary}
\newenvironment{proof}[1][Proof]{\noindent\textbf{#1.} }{\ \rule{0.5em}{0.5em}}
\input{tcilatex}
\begin{document}


\begin{center}
\begin{equation*}
\FRAME{itbpF}{1.0525in}{0.7515in}{0in}{}{}{Figure}{\special{language
"Scientific Word";type "GRAPHIC";maintain-aspect-ratio TRUE;display
"USEDEF";valid_file "T";width 1.0525in;height 0.7515in;depth
0in;original-width 2.3168in;original-height 1.6466in;cropleft "0";croptop
"1";cropright "1";cropbottom "0";tempfilename
'TCTQIQ04.bmp';tempfile-properties "XPR";}}
\end{equation*}%
\textbf{F\'{\i}sica II (2026-1)}

\textbf{Tarea 1}\emph{: Fuerza y Campo el\'{e}ctrico de cargas puntuales y
distribuciones continuas}

\textit{Profesor: Arturo Fern\'{a}ndez P\'{e}rez}

\textbf{Plazo de Entrega: Mi\'{e}rcoles 16 de abril a las 21:00 hrs.}
\end{center}

\begin{enumerate}
\item Tres cargas puntuales $q_{1}=11,8$ $[nC]$, $q_{2}=-31,5$ $[nC]$ y $%
q_{3}=-10,3$ $[nC]$\ se ubican en un sistema de referencia, cuyas
coordenadas y posiciones (en cent\'{\i}metros, [cm]) se muestran en la Fig.
1. La carga $q_{1}$ se encuentra a $8,0$ [cm] del origen. Determine:

\begin{enumerate}
\item La fuerza el\'{e}ctrica neta $\vec{F}$ ejercida sobre la carga $q_{2}.$

\item El campo el\'{e}ctrico neto $\vec{E}$ en el punto (1,0) [cm]%
\begin{equation*}
\FRAME{itbpFU}{2.0427in}{2.3886in}{0in}{\Qcb{Figura 1}}{}{Figure}{\special%
{language "Scientific Word";type "GRAPHIC";maintain-aspect-ratio
TRUE;display "USEDEF";valid_file "T";width 2.0427in;height 2.3886in;depth
0in;original-width 7.4149in;original-height 8.6836in;cropleft "0";croptop
"1";cropright "1";cropbottom "0";tempfilename
'TCTPZ400.bmp';tempfile-properties "XPR";}}
\end{equation*}
\end{enumerate}

\item Dada la siguiente configuraci\'{o}n de la Fig. 2, con $q_{1}=q_{2}=2Q$
y $q_{3}=-Q$, determinar la fuerza neta que experimenta una part\'{\i}cula $%
q_{0}=Q$ ubicada en el punto P. Considere $Q=10$ [nC] y $R=25$ [cm]. Dibuje
las l\'{\i}neas de fuerza que cada carga induce sobre la carga $q_{0}$%
\begin{equation*}
\FRAME{itbpFU}{1.8844in}{1.8023in}{0in}{\Qcb{Figura 2}}{}{Figure}{\special%
{language "Scientific Word";type "GRAPHIC";maintain-aspect-ratio
TRUE;display "USEDEF";valid_file "T";width 1.8844in;height 1.8023in;depth
0in;original-width 2.3644in;original-height 2.2606in;cropleft "0";croptop
"1";cropright "1";cropbottom "0";tempfilename
'TCTQ5G01.wmf';tempfile-properties "XPR";}}
\end{equation*}%
\pagebreak 

\item Una barra con densidad lineal de carga uniforme $\lambda =-3,6$ [$nC/m$%
] se ubica sobre el eje $x$, como muestra la Fig. 3, entre $x=5$ [cm] e $x=18
$ [cm\.{]}. Adem\'{a}s, en el punto (-2,6) [cm] se ubica una carga puntual $%
q=-16$ [$nC$]. Determine:

\begin{enumerate}
\item El campo el\'{e}ctrico $\vec{E}$ inducido por la barra en el punto
(-2,6) [cm].

\item La fuerza el\'{e}ctrica $\vec{F}$ sobre la carga $q.$%
\begin{equation*}
\FRAME{itbpFU}{3.8527in}{2.1041in}{0in}{\Qcb{Figura 3}}{}{Figure}{\special%
{language "Scientific Word";type "GRAPHIC";maintain-aspect-ratio
TRUE;display "USEDEF";valid_file "T";width 3.8527in;height 2.1041in;depth
0in;original-width 12.3167in;original-height 6.7075in;cropleft "0";croptop
"1";cropright "1";cropbottom "0";tempfilename
'TCTQCF02.bmp';tempfile-properties "XPR";}}
\end{equation*}
\end{enumerate}

\item Considere un \emph{semianillo }(la mitad de una circunferencia) el cu%
\'{a}l posee una densidad de carga lineal uniforme $\lambda =18,0$ [$\mu C/m$%
] y un radio $R=30,0$ [cm] como muestra la Fig. 4.

\begin{itemize}
\item Determine el campo el\'{e}ctrico $\vec{E}$ en el eje del \emph{%
semianillo}, a una distancia $z=21,0$ [cm] de su centro.

\item Si en ese punto se ubica una carga puntual $q=22,0$ [$\mu C]$
determine la fuerza el\'{e}ctrica $\vec{F}$ ejercida sobre ella.%
\begin{equation*}
\FRAME{itbpFU}{2.405in}{2.1776in}{0in}{\Qcb{Figura 2}}{}{Figure}{\special%
{language "Scientific Word";type "GRAPHIC";maintain-aspect-ratio
TRUE;display "USEDEF";valid_file "T";width 2.405in;height 2.1776in;depth
0in;original-width 4.6164in;original-height 4.1753in;cropleft "0";croptop
"1";cropright "1";cropbottom "0";tempfilename '../clases
I-2021/230045/TCTQHM03.wmf';tempfile-properties "XPR";}}
\end{equation*}
\end{itemize}
\end{enumerate}

\end{document}
